\documentclass[conference]{IEEEtran}

\usepackage{cite}
\usepackage{amsmath,amssymb,amsfonts}
\usepackage{algorithmic}
\usepackage{graphicx}
\usepackage{textcomp}
\usepackage{xcolor}
\usepackage{booktabs}
\usepackage{url}
\usepackage{tagging}

\usepackage{sc26repro}

%\usetag{explanation}
%\usetag{example}

\begin{document}

\twocolumn[%
{\begin{center}
\Huge
Appendix: Artifact Description/Artifact Evaluation
\end{center}}
]

%%%%%%%%%%%%%%%%%%%%%%%%%%%%%%%%%%%%%%%%%%%%%%%%%%%%%
%  AD Appendix
%%%%%%%%%%%%%%%%%%%%%%%%%%%%%%%%%%%%%%%%%%%%%%%%%%%%%

\appendixAD

\section{Overview of Contributions and Artifacts}

\subsection{Paper's Main Contributions}

\begin{description}
\item[$C_1$] \textbf{Bid-based reclamation interface.}
  A structured per-request bid abstraction that surfaces recoverable capacity
  and estimated disruption cost to the runtime scheduler, replacing implicit
  victim-ordering rules with an explicit decision interface.

\item[$C_2$] \textbf{Decoupled scoring-to-selection architecture.}
  Separation of request-level disruption estimation from coordinated
  cross-request victim selection through a uniform cost--capacity bid
  interface, so that the selection algorithm is independent of how
  disruption estimates are produced.

\item[$C_3$] \textbf{Portable runtime integration.}
  A non-invasive scheduling layer with a \texttt{FrameworkAdapter} abstraction
  demonstrated on vLLM~0.17.1 and SGLang, reusing each framework's native
  reclamation and recovery mechanisms without source-code modification.

\item[$C_4$] \textbf{Structured evaluation under controlled conditions.}
  Evaluation of five strategies varying across four scheduling dimensions
  (admission ordering, running-queue reordering, victim selection, proactive
  reclamation) under frozen request traces and calibrated arrival rates,
  enabling structured attribution across these co-varying dimensions.
\end{description}


\subsection{Computational Artifacts}

\begin{description}
\item[$A_1$] BidKV source code and experiment framework:\\
  \url{https://github.com/[anonymized]/bidkv}
\end{description}

\begin{center}
\resizebox{\columnwidth}{!}{%
\begin{tabular}{@{}rll@{}}
\toprule
Artifact ID & Contributions & Related Paper Elements \\
            & Supported     & \\
\midrule
$A_1$  & $C_1, C_2, C_3, C_4$ & Table~I (strategy differentiation) \\
       &                      & Table~II (rate sensitivity, vLLM) \\
       &                      & Table~III (long-context results) \\
       &                      & Table~IV (SGLang portability) \\
       &                      & Table~V (robustness sensitivity) \\
       &                      & Figure~1 (motivation) \\
       &                      & Figure~2 (system architecture) \\
       &                      & Figure~3 (algorithm) \\
       &                      & Figure~4 (reclamation analysis) \\
       &                      & Figure~5 (victim-selection tradeoff) \\
\bottomrule
\end{tabular}%
}
\end{center}


%%%%%%%%%%%%%%%%%%%%%%%%%%%%%%%%%%%%%%%%%%%%%%%%%%%%%%%%
\section{Artifact Identification}
%%%%%%%%%%%%%%%%%%%%%%%%%%%%%%%%%%%%%%%%%%%%%%%%%%%%%%%%

\newartifact

\artrel

$A_1$ comprises (i)~the BidKV Python package
(\texttt{src/bidkv/}), (ii)~vLLM and SGLang adapter layers
(\texttt{src/bidkv/adapters/}), (iii)~five baseline strategy
implementations (\texttt{src/bidkv/baselines/}), and (iv)~the
experiment runner, trace-freeze utilities, and analysis scripts
(\texttt{src/bidkv/experiments/}).
Together these components instantiate $C_1$--$C_3$ and drive all
measurements described in $C_4$.

\artexp

Among all five vLLM strategies at rate~3.8\,req/s (mixed workload), BidKV
should achieve the best P95 TTFT ($\approx$631\,ms) and near-best 300\,ms
SLO attainment ($\approx$83.2\,\%, second only to Static-Random at 84.1\,\%)
(Table~II).
Across all three arrival rates (2.0, 3.8, and 5.7\,req/s), BidKV should
rank first or second on both P95 TTFT and SLO attainment (Table~II, Figure~4).
On the long-context workload at peak load (rate~0.70\,req/s), BidKV should
achieve the lowest P95 TTFT among all five strategies
($\approx$30{,}143\,ms; Table~III).
SGLang portability results (Table~IV) are pending and will be included in the
camera-ready version.
These results substantiate $C_1$--$C_4$: the bid interface and
utility-ranked selection produce measurably better admission responsiveness
across two arrival-rate regimes and two workload types.

\arttime

\begin{itemize}
  \item \textbf{Artifact Setup:} 20--30\,min
    (conda environment creation, package installation, model download if not
    cached, trace generation).
  \item \textbf{Artifact Execution:}
    vLLM full matrix (5 strategies $\times$ 2 workloads $\times$ 3 rates
    $\times$ 3 runs = 90 experiment runs): $\approx$ 20--28\,h
    (mixed-workload runs average 5--8\,min each; long-context runs average
    20--25\,min each; server startup adds $\approx$4\,min per run).
    SGLang portability slice: timing pending (results not yet included in
    this artifact version).
  \item \textbf{Artifact Analysis:} 5--10\,min
    (single analysis script invocation per result directory).
\end{itemize}

\artin

\artinpart{Hardware}

\begin{itemize}
  \item Single NVIDIA RTX A6000 GPU (48\,GiB VRAM); CUDA~12.5.
  \item $\geq$64\,GiB system RAM; $\geq$8 CPU cores recommended.
  \item $\geq$40\,GB free disk space (model weights, result files).
  \item No special interconnect is required; single-node, single-GPU only.
\end{itemize}

\artinpart{Software}

\begin{itemize}
  \item Python~3.10 or 3.11
  \item conda (Miniconda/Anaconda) for environment management
  \item vLLM~0.17.1 (\url{https://github.com/vllm-project/vllm})
  \item SGLang (version matching \texttt{requirements.txt};\\
    \url{https://github.com/sgl-project/sglang})
  \item BidKV package (this artifact; zero external dependencies beyond
    vLLM/SGLang, which are runtime-only imports in the adapter layer)
\end{itemize}

\artinpart{Datasets / Inputs}

The experiment uses the ShareGPT Vicuna unfiltered dataset
(\url{https://huggingface.co/datasets/anon8231489123/}).
Request traces are deterministically generated from this dataset using
seed~42 and frozen by SHA-256 checksums.
The trace-freeze script is included in the artifact:

\begingroup\footnotesize\begin{verbatim}
python -m bidkv.experiments.vllm.freeze_traces \
    --mode formal --use-frozen-rates \
    --output-dir experiments/vllm/traces/
\end{verbatim}
\endgroup

The model weights are Llama-3.1-8B-Instruct
(\url{https://huggingface.co/meta-llama/Llama-3.1-8B-Instruct}).
If a local HuggingFace cache is available, set
\texttt{HF\_HUB\_OFFLINE=1} and \texttt{TRANSFORMERS\_OFFLINE=1}.

\artinpart{Installation and Deployment}

\begingroup\footnotesize\begin{verbatim}
# 1. Create conda environment
conda create -n bidkv python=3.10 -y
conda activate bidkv

# 2. Install vLLM and SGLang
pip install vllm==0.17.1
pip install sglang[all]

# 3. Install BidKV (editable)
cd /path/to/bidkv
pip install -e ".[dev]"

# 4. Generate frozen traces (once)
python -m bidkv.experiments.vllm.freeze_traces \
    --mode formal --use-frozen-rates \
    --output-dir experiments/vllm/traces/
\end{verbatim}
\endgroup

No custom compiler is required.
All components are pure Python.

\artcomp

\textbf{Workflow tasks and dependencies:}

\begin{itemize}
  \item $T_1$ \textbf{Trace generation.}
    \texttt{freeze\_traces.py} samples 1\,000 requests (mixed) and 500
    requests (long-context) from ShareGPT with seed~42 and writes
    SHA-256-verified JSON trace files to \texttt{experiments/vllm/traces/}.

  \item $T_2$ \textbf{vLLM experiment execution.}
    \texttt{runner.py} iterates over 5 strategies $\times$ 2 workloads
    $\times$ 3 arrival rates $\times$ 3 runs.
    For each strategy it starts a vLLM server (with BidKV adapter injected
    via plugin for non-PE strategies), replays the frozen trace using a
    Poisson open-loop arrival model, collects per-request TTFT and
    completion-token counts, and writes one JSON result file per run.
    Fixed server parameters: \texttt{--gpu-memory-utilization~0.5},
    \texttt{--num-gpu-blocks-override~600}, \texttt{--max-num-seqs~32},
    \texttt{--block-size~16}, \texttt{--max-model-len~8192},
    \texttt{--enforce-eager}, \texttt{--disable-frontend-multiprocessing},
    \texttt{--no-enable-prefix-caching}.

  \item $T_3$ \textbf{SGLang portability execution (pending).}
    Uses the SGLang runner (\texttt{experiments/sglang/runner.py}) with
    3 strategies (BidKV, Slack-Aware, vanilla SGLang) $\times$
    rate~3.8\,req/s $\times$ 3 runs.
    Results for Table~IV are not yet included in this artifact version.

  \item $T_4$ \textbf{Analysis.}
    \texttt{analysis.py} reads all JSON result files, computes P95 TTFT,
    TPOT P95, throughput, and SLO attainment (from raw per-request data;
    P95 and SLO are not pre-computed), and produces tables and figures.

\end{itemize}

\textbf{Dependencies:} $T_1 \rightarrow T_2 \rightarrow T_4$;
$T_1 \rightarrow T_3 \rightarrow T_4$.
$T_2$ and $T_3$ are independent and can run in parallel if GPUs are available.

\textbf{Experimental parameters:}
\begin{itemize}
  \item Mixed workload: 1\,000 requests, rates $\{2.0, 3.8, 5.7\}$\,req/s,
    3 runs per (strategy, rate) combination.
  \item Long-context workload: 500 requests,
    rates $\{0.35, 0.5, 0.7\}$\,req/s, 3 runs per combination.
  \item SGLang portability ($T_3$, pending): mixed workload,
    rate~3.8\,req/s, 3 runs per strategy.
  \item KV capacity fixed at 600 blocks (9\,600 tokens) to induce
    sustained KV pressure; this is the critical parameter for
    strategy differentiation and must not be changed.
  \item Arrival process: Poisson (open-loop); seed~42 for trace replay.
\end{itemize}

\artout

Raw results are stored as JSON files with naming convention
\texttt{strategy\_\_workload\_\_rate\{r\}\_\_r\{i\}.json}
in the designated result directory.
Each file contains per-request TTFT, total latency, completion-token count,
and adapter-level reclamation metrics.

Analysis scripts in \texttt{experiments/vllm/} and
\texttt{experiments/sglang/} read these files and compute P95/P50
TTFT, TPOT P95, throughput, and SLO attainment from raw
per-request data (P95 and SLO are not pre-computed in the JSON).
Outputs are CSV tables (Table~II--III) and PDF figures
(Figures~4--5).
The scripts handle the legacy field-name difference between older
result batches (\texttt{total\_compressions}) and newer ones
(\texttt{total\_evictions}).

\end{document}
